% !PRINTED 99
\weberpage{99}

Aus der Gleichung

\weq{8}{\alpha\delta - \beta\gamma = 1}

folgt, daß weder $\alpha, \beta$ noch $\alpha, \gamma$, noch $\delta, \beta$, noch $\delta, \gamma$ einen gemeinschaftlichen Faktor haben können. Hat man aber $\alpha, \beta$ beliebig ohne gemeinschaftlichen Teiler angenommen, so kann man $\gamma, \delta$ noch auf unendlich viele Arten aus (8) bestimmen. Ist $\gamma, \delta$ eine dieser Bestimmungen, so sind sie alle in der Form

\weq{9}{\gamma + \lambda\alpha, \quad \delta + \lambda\beta}

enthalten, worin $\lambda$ eine beliebige ganze Zahl ist (Bd. I, \S\ 126).

\begin{center}
\textbf{\S\ 29.\quad\gesperrt{Zusammensetzung der Transformationen aus einfacheren.}}
\end{center}

In dem System $\mathfrak{S}$ aller Transformationen $S$ ist ein System $\mathfrak{S}_0$ enthalten, das aus allen den Transformationen $S_0$ besteht, deren zweite Transformationszahl $b = 0$ ist, während $a$ und $\partial$ \gesperrt{positiv} sind:

\weq{1}{S_0 = \begin{pmatrix} a, & 0 \\ c, & \partial \end{pmatrix}.}

Bei der Zusammensetzung zweier $S_0$ entsteht wieder ein $S_0$, aber doch ist $\mathfrak{S}_0$ so wenig eine Gruppe wie $\mathfrak{S}$.

Man kann jede beliebige Transformation
\[
S = \begin{pmatrix} p, & q \\ r, & s \end{pmatrix}
\]
durch eine Zusammensetzung $LS$ auf ein $S_0$ zurückführen. Soll nämlich

\weq{2}{\begin{pmatrix} \alpha, & \beta \\ \gamma, & \delta \end{pmatrix}\begin{pmatrix} p, & q \\ r, & s \end{pmatrix} = \begin{pmatrix} a, & 0 \\ c, & \partial \end{pmatrix}}

sein, so muß $\alpha, \beta$ der Bedingung genügen:
\[
\alpha q + \beta s = 0,
\]
und wenn also $\partial$ der größte gemeinschaftliche Teiler von $q$ und $s$ ist, so setze man
\[
\partial\alpha = s, \quad \partial\beta = -q,
\]
und bestimme, nachdem $\alpha$ und $\beta$ hierdurch als relative Primzahlen ermittelt sind, $\gamma$ und $\delta$ aus der Formel (8), \S\ 28. Dann ist (2) erfüllt, wenn
\[
a\partial = n, \quad c = \gamma p + \delta r
\]
gesetzt wird; $a$ und $\partial$ sind hierdurch eindeutig bestimmt, $c$ kann aber bei anderer Wahl von $\gamma$ und $\delta$ durch $c+\lambda a$ ersetzt werden. Man kann daher über $\lambda$ so verfügen, daß $c$ in der Reihe
%CONT>>
